\scanpage{40}\mspage{39}
pour $ \tau'$. Donc cela impose, comme parties fermées, pour $\tau'$, déjà au moins les
\[
\left\{\begin{aligned}
 A\subset S&\quad\text{fermés pour }\tau\\
 B\subset T&\quad\text{fermés pour }\tau'_0,
\end{aligned}\right.
\]
donc les réunions $A\cup B$. Je dis que ceux-ci sont \inserted{déjà} les fermés d'une ultr.topologie (qui visiblement raffine $\tau$, et qui induit la topologie $\tau'_0$, puisque
\[
 (A\cup B)\cap T=(A\cap T)\cup B
\]
et $A\cap T$ est fermé pour $\tau'_0$, donc on trouve exactement les fermés de $(T,\tau'_0)$).
% slanted
\margin{façon \uline{canonique} d'écrire un tel fermé $X$ de $S$ en termes de $\deleted{\unsure{A}}B=X\cap T$ et $U=X\cap(S\setminus T)$, $U$ un \uline{fermé de $S\setminus T$} (pour $\tau\mid S\setminus T$) condition $\boxed{B\supset\overline U\inserted{\text{adh. pour }\tau\text{, bien sûr}}\cap T}$ et on a $X=B\cup\overline U$. Cette fois correspondance biunivoque entre couples $(B,U)$, satisfaisant les conditions dites, et les $X$. (Dans la représentation précédente, on prend $A=\overline U$, donc on suppose $1^\circ)\ B\supset A\cap T\quad 2^\circ)\ A\cap(S\setminus T)$ dense \unsure{en|dans} $A$.}
\unsure{ce sera|c'est} donc la \deleted{plus} moins fine des topologies cherchées.

\deleted{\unsure{Déjà}} Stabilité par sommes quelconques : triviale. Stabilité par intersections quelconques
\[
 \bigcap_{i\in I}(A_i\cup B_i)=
 \bigcup_{J\in I}
 \biggl(\underset{\underset{A}{\rule{.35pt}{.3em}\rule{3em}{.35pt}\rule{.35pt}{.3em}}}{\bigcap_{i\in J}A_i\biggr)}
 \cap\biggl(\underset{\underset{B}{\rule{.35pt}{.3em}\rule{2.8em}{.35pt}\rule{.35pt}{.3em}}}{\bigcap_{\unsure{j|J}\in I\setminus J} B_J}\biggr)
\]
c'est une réunion de fermés de $(S,\tau)$, donc un fermé de $(S,\tau)$, OK.

Quelle est la topologie induite \inserted{par $\tau'$} sur $S\setminus T$ ? $(A\cup B)\cap(S\setminus T)=A\cap(S\setminus T)$, on trouve la top. induite par $\tau$.

\deleted{Inversement, si $\tau'\leq\tau$, $\tau'\mid T=\tau'_0$, et $\tau'\mid S\setminus T=\tau\mid S\setminus\deleted{\unsure{J}}T$, je dis que $\tau'$ est la top. précédente. \unsure{Il reste à|Je voudrais} voir que ses fermés sont de la forme $A\cup B$. Mais soit un tel fermé $C$, alors \deleted{\unsure{$C=C$}}}
